\chapter{การจำลองครึ่งชีวิตกัมมันตภาพรังสีและประสิทธิภาพรังสี UV-C ฆ่าเชื้อ}

\begin{planbox}[สัปดาห์ที่ 12 — เวลาปฏิบัติการ 3 ชั่วโมง (สอดคล้องกับ CLO1, CLO2 และ CLO4)]
\textbf{1. วัตถุประสงค์การทดลอง}
\begin{itemize}
    \item ทำการทดลองจำลองการสลายตัวของนิวเคลียสกัมมันตรังสีด้วยชุดลูกเต๋าหรือเหรียญสองหน้า เพื่อพิสูจน์สมการครึ่งชีวิต $N = N_0(1/2)^n$ โดยไม่ต้องใช้แคลคูลัส
    \item วัดความเข้มรังสีตามระยะห่างจากแหล่งกำเนิดแสงเพื่อพิสูจน์กฎกำลังสองผกผัน ($I \propto 1/r^2$)
    \item วิเคราะห์ข้อมูลการลดลงของจำนวนโคโลนีแบคทีเรียเมื่อได้รับรังสี UV-C ($254\text{ nm}$) ที่ระยะเวลาและระยะห่างต่างกัน เพื่อหาค่าเวลาลดทอนครึ่งชีวิตของประชากรเชื้อ ($D_{1/2}$)
\end{itemize}
\textbf{2. ความเชื่อมโยงกับผลลัพธ์การเรียนรู้ของรายวิชา (CLOs) และหลักสูตร (PLOs)}
\begin{itemize}
    \item \textbf{CLO1 (PLO2)}\quad อธิบายแนวคิดครึ่งชีวิต (Half-life) และกฎกำลังสองผกผันของความเข้มรังสีได้
    \item \textbf{CLO2 (PLO2, PLO3)}\quad ประยุกต์ใช้สมการครึ่งชีวิตและกฎกำลังสองผกผันในการวางแผนระยะห่างและเวลาเปิดหลอดรังสี UV-C ฆ่าเชื้อในตู้ปลอดเชื้อ (Biosafety Cabinet) ได้
    \item \textbf{CLO4 (PLO5, PLO6)}\quad รวบรวมข้อมูลสถิติของกลุ่ม เขียนกราฟการสลายตัว และสื่อสารผลการทดลองได้อย่างถูกต้อง
\end{itemize}
\end{planbox}

\section{หลักการทฤษฎีและการเชื่อมโยงสู่งานจุลชีววิทยา}

การสลายตัวของสารกัมมันตรังสีและการลดลงของจำนวนจุลินทรีย์ที่มีชีวิตรอดเมื่อได้รับรังสีฆ่าเชื้อ มีรูปแบบทางคณิตศาสตร์เดียวกัน คือ \textbf{การลดลงแบบครึ่งชีวิต} (\textbf{Half-Life Decay})\index{ครึ่งชีวิต (Half-Life)} โดยในทุกๆ ช่วงเวลา 1 ครึ่งชีวิต ($T_{1/2}$) ปริมาณสารตั้งต้นหรือจำนวนจุลินทรีย์จะลดลงเหลือ\textbf{ครึ่งหนึ่ง}ของปริมาณที่มีอยู่เดิมเสมอ หากผ่านไป $n$ รอบครึ่งชีวิต ($n = t / T_{1/2}$) ปริมาณที่เหลืออยู่ $N$ คำนวณได้จากสมการที่ \eqref{eq:lab11_halflife}

\begin{equation}
    N = N_0 \left(\frac{1}{2}\right)^n \qquad \text{เมื่อ} \qquad n = \frac{t}{T_{1/2}}
    \label{eq:lab11_halflife}
\end{equation}

สำหรับการฆ่าเชื้อบนพื้นผิวด้วยหลอดรังสีอัลตราไวโอเลตชนิด \textbf{UV-C} ($\lambda = 254\text{ nm}$)\index{รังสีอัลตราไวโอเลต (UV-C Sterilization)} ซึ่งทำลายพันธุกรรมของแบคทีเรียด้วยการสร้างพันธะไทมีนไดเมอร์ (Thymine Dimer) บนสายดีเอ็นเอ ความเข้มรังสี $I$ (หน่วย $\text{W/m}^2$ หรือ $\mu\text{W/cm}^2$) ที่ตกกระทบบนจานเพาะเชื้อจะลดลงอย่างรวดเร็วเมื่อระยะห่าง $r$ จากหลอดไฟเพิ่มขึ้น ตาม\textbf{กฎกำลังสองผกผัน} (\textbf{Inverse-Square Law})\index{กฎกำลังสองผกผัน (Inverse-Square Law)}

\begin{conceptbox}[สรุปสูตรคำนวณฉบับเข้าใจง่ายสำหรับปฏิบัติการที่ 11]
\begin{itemize}
    \item \textbf{สมการครึ่งชีวิตแบบคิดเลขในใจ}\quad $N = \dfrac{N_0}{2^n}$ (เช่น ผ่านไป $3$ ครึ่งชีวิต เหลือ $N_0 / 2^3 = N_0 / 8$)
    \item \textbf{กฎกำลังสองผกผันของความเข้มรังสีเมื่อเทียบระยะห่าง $r_1$ กับ $r_2$}
    \begin{equation*}
        \frac{I_2}{I_1} = \left(\frac{r_1}{r_2}\right)^2 \implies I_2 = I_1 \left(\frac{r_1}{r_2}\right)^2
    \end{equation*}
    \item \textbf{ปริมาณรังสีสะสม (UV Dose)}\quad $\text{Dose} = \text{ความเข้มรังสี } I \times \text{เวลาฉายรังสี } t$
\end{itemize}
\end{conceptbox}

\begin{worked}[การคำนวณความเข้มรังสี UV-C และจำนวนแบคทีเรียที่รอดชีวิต]
\textbf{โจทย์} หลอดรังสี UV-C ในตู้ปลอดเชื้อให้ความเข้มรังสี $I_1 = 400\ \mu\text{W/cm}^2$ ที่ระยะห่าง $r_1 = 20\text{ cm}$ หากวางจานเพาะเชื้อห่างออกไปเป็น $r_2 = 40\text{ cm}$ ความเข้มรังสี $I_2$ จะเหลือเท่าใด และหากที่ระยะนี้แบคทีเรียมีเวลาครึ่งชีวิตการตาย $T_{1/2} = 2\text{ นาที}$ เมื่อเปิดไฟนาน $t = 8\text{ นาที}$ จากเชื้อเริ่มต้น $N_0 = 1{,}600\text{ CFU}$ จะเหลือรอดกี่ $\text{CFU}$

\textbf{วิธีทำทีละขั้นตอน}
\begin{enumerate}
    \item คำนวณความเข้มรังสี $I_2$ ที่ระยะ $r_2 = 40\text{ cm}$ (ห่างขึ้น $2$ เท่า ความเข้มลดลง $2^2 = 4$ เท่า)
    \begin{equation*}
        I_2 = 400 \times \left(\frac{20}{40}\right)^2 = 400 \times \frac{1}{4} = 100\ \mu\text{W/cm}^2
    \end{equation*}
    \item คำนวณจำนวนรอบครึ่งชีวิต $n = t / T_{1/2} = 8 / 2 = 4\text{ รอบ}$ แล้วหาจำนวนเชื้อที่เหลือรอด $N$
    \begin{equation*}
        N = \frac{1{,}600}{2^4} = \frac{1{,}600}{16} = 100\text{ CFU}
    \end{equation*}
\end{enumerate}
\textbf{คำตอบ} ความเข้มรังสีเหลือ $100\ \mu\text{W/cm}^2$ และมีแบคทีเรียเหลือรอดชีวิต $100\text{ CFU}$
\end{worked}

\begin{figure}[H]
\centering
\resizebox{0.86\textwidth}{!}{%
\begin{tikzpicture}[>=Stealth, font=\small]
    % Left: Half-life decay curve
    \draw[->, line width=1.1pt] (0,0) -- (6.2,0) node[right, font=\footnotesize\bfseries] {จำนวนรอบครึ่งชีวิต $n$};
    \draw[->, line width=1.1pt] (0,0) -- (0,3.8) node[above, font=\footnotesize\bfseries] {จำนวนที่เหลือ $N/N_0$ (\%)};
    \draw[line width=1.8pt, physRoyalBlue, domain=0:5, samples=60] plot (\x, {3.2*(0.5^\x)});
    \foreach \n/\pct/\y in {0/100\%/3.2, 1/50\%/1.6, 2/25\%/0.8, 3/12.5\%/0.4, 4/6.25\%/0.2} {
        \fill[crimsonAlert] (\n,\y) circle (2.5pt);
        \draw[dashed, gray] (\n,0) -- (\n,\y) -- (0,\y);
        \node[below, font=\scriptsize] at (\n,0) {$\n$};
        \node[left, font=\scriptsize] at (0,\y) {\pct};
    }

    % Right: UV-C Inverse-square law in Biosafety Cabinet
    \draw[fill=purple!25, draw=purple!80!black, line width=1.2pt, rounded corners=3pt] (8.5,3.5) rectangle (12.5,3.85);
    \node[font=\footnotesize\bfseries\color{purple!80!black}, above] at (10.5,3.85) {หลอดรังสีฆ่าเชื้อ UV-C ($254\text{ nm}$)};
    % Rays
    \draw[->, line width=1.1pt, purple!75] (10.0,3.5) -- (8.2,0.5);
    \draw[->, line width=1.1pt, purple!75] (10.5,3.5) -- (10.5,0.5);
    \draw[->, line width=1.1pt, purple!75] (11.0,3.5) -- (12.8,0.5);

    % Distance r1
    \draw[fill=boxBlueBg, draw=physRoyalBlue, line width=1pt] (9.4,2.1) rectangle (11.6,2.35);
    \node[right, font=\scriptsize\bfseries\color{physRoyalBlue}] at (11.7,2.22) {ระยะ $r_1 = 20\text{ cm}$ ($I_1 = 100\%$)};

    % Distance r2 = 2 r1
    \draw[fill=boxAmberBg, draw=physAmber, line width=1pt] (8.3,0.7) rectangle (12.7,0.95);
    \node[right, font=\scriptsize\bfseries\color{physAmber!85!black}] at (12.8,0.82) {ระยะ $r_2 = 40\text{ cm}$ ($I_2 = 25\%$)};
\end{tikzpicture}%
}
\caption{กราฟการลดลงแบบครึ่งชีวิต (ซ้าย) และการลดลงของความเข้มรังสี UV-C ตามกฎกำลังสองผกผันของระยะห่าง (ขวา)}
\label{fig:lab11_setup}
\end{figure}

\section{วัสดุอุปกรณ์และขั้นตอนการปฏิบัติการ}

\begin{equipbox}
\begin{enumerate}
    \item ชุดเบี้ยสองสีหรือเหรียญจำลองนิวเคลียสกัมมันตรังสี จำนวน $100$ ชิ้น พร้อมกล่องเขย่า จำนวน $1$ ชุด
    \item แหล่งกำเนิดแสงจำลองหลอด UV และเครื่องวัดความเข้มแสง (Lux Meter / UV Radiometer) พร้อมรางวัดระยะทาง
    \item ชุดภาพถ่ายจานเพาะเชื้อแบคทีเรีย \textit{Escherichia coli} ที่ผ่านการฉายรังสี UV-C ($254\text{ nm}$) เป็นเวลา $0, 1, 2, 3$ และ $4$ นาที สำหรับฝึกนับจำนวนโคโลนี (CFU)
\end{enumerate}
\end{equipbox}

\begin{warningbox}
ในการปฏิบัติงานจริงกับตู้ปลอดเชื้อที่มีหลอดรังสี UV-C ต้องปิดกระจกหน้าตู้ลงให้สนิทและห้ามสอดมือหรือจ้องมองแสงสีม่วงฟ้าของหลอด UV-C โดยตรง เพราะจะทำให้เกิดกระจกตาอักเสบและผิวหนังไหม้ได้
\end{warningbox}

\subsection{ขั้นตอนการทดลอง}
\begin{enumerate}
    \item \textbf{ตอนที่ 1 การจำลองการสลายตัวครึ่งชีวิตด้วยเบี้ยสองสี}\quad เริ่มต้นด้วยเบี้ยจำนวน $N_0 = 100$ ชิ้น เขย่ากล่องแล้วเทลงบนถาด คัดเบี้ยที่หงายด้านสีแดง (แทนนิวเคลียสที่สลายตัวแล้วหรือเซลล์ที่ตายแล้ว) ออก นับจำนวนเบี้ยที่เหลือรอด $N$ บันทึกเป็นรอบการเขย่าที่ $n = 1, 2, 3, 4$ และ $5$ เปรียบเทียบกับค่าทางทฤษฎีใน\mbox{ตารางที่} 11.1
    \item \textbf{ตอนที่ 2 การทดสอบกฎกำลังสองผกผันและการนับโคโลนีรอดชีวิตจากรังสี UV-C}\quad วัดความเข้มแสงที่ระยะห่าง $r = 15, 30$ และ $45\text{ cm}$ จากหลอดกำเนิดแสง และนับจำนวนโคโลนีแบคทีเรียที่รอดชีวิตบนจานเพาะเชื้อตามเวลาฉายรังสี UV-C บันทึกใน\mbox{ตารางที่} 11.2
\end{enumerate}

\section{ใบบันทึกผลการทดลองและคำถามท้ายปฏิบัติการ}

\begin{worksheetbox}[ใบบันทึกผลการทดลองที่ 11\quad ครึ่งชีวิตและประสิทธิภาพรังสี UV-C ฆ่าเชื้อ]
\textbf{กลุ่มปฏิบัติการที่} \dotfill \textbf{วันที่ทำการทดลอง} \dotfill\\[4pt]
\textbf{รายชื่อสมาชิกกลุ่ม} 1. \dotfill 2. \dotfill 3. \dotfill

\vspace{4pt}
\textbf{ตารางที่ 11.1\quad ผลการจำลองการลดลงแบบครึ่งชีวิตจากจำนวนเริ่มต้น $N_0 = 100$ ชิ้น}
\begin{center}
\footnotesize\setlength{\tabcolsep}{3.5pt}
\begin{tabularx}{\linewidth}{c c c c >{\raggedright\arraybackslash}X}
\toprule
\textbf{รอบที่ ($n$)} & \textbf{ทฤษฎี ($100/2^n$)} & \textbf{เหลือจริง (1)} & \textbf{เหลือจริง (2)} & \textbf{ค่าเฉลี่ยจากการทดลอง $\bar{N}$} \\
\midrule
$n = 0$ (เริ่มต้น) & $100.0$ & $100$ & $100$ & $100.0$ \\[4pt]
$n = 1$ & $50.0$ & \dotfill & \dotfill & \dotfill \\[4pt]
$n = 2$ & $25.0$ & \dotfill & \dotfill & \dotfill \\[4pt]
$n = 3$ & $12.5$ & \dotfill & \dotfill & \dotfill \\[4pt]
$n = 4$ & $6.25$ & \dotfill & \dotfill & \dotfill \\
\bottomrule
\end{tabularx}
\end{center}

\vspace{4pt}
\textbf{ตารางที่ 11.2\quad ผลการวัดความเข้มรังสีตามระยะห่างและจำนวนโคโลนีที่รอดชีวิตจากรังสี UV-C}
\begin{center}
\footnotesize\setlength{\tabcolsep}{3.5pt}
\begin{tabularx}{\linewidth}{c c c | c c >{\raggedright\arraybackslash}X}
\toprule
\textbf{ระยะ $r$} & \textbf{ความเข้ม $I$} & \textbf{$I/I_{15}$} & \textbf{เวลา UV} & \textbf{โคโลนี (CFU)} & \textbf{ร้อยละการรอดชีวิต (\%)} \\
\midrule
$15\text{ cm}$ ($1\times$) & ........ & $1.00$ & $0\text{ min}$ & $400\text{ CFU}$ & $100\%$ \\[6pt]
$30\text{ cm}$ ($2\times$) & ........ & ........ & $1\text{ min}$ & ........ & \dotfill \\[6pt]
$45\text{ cm}$ ($3\times$) & ........ & ........ & $2\text{ min}$ & ........ & \dotfill \\[6pt]
$60\text{ cm}$ ($4\times$) & ........ & ........ & $4\text{ min}$ & ........ & \dotfill \\
\bottomrule
\end{tabularx}
\end{center}

\vspace{4pt}
\textbf{สรุปและอภิปรายผลการทดลอง}\\[4pt]
\dotfill\\[6pt]
\dotfill
\end{worksheetbox}

\subsection{คำถามท้ายปฏิบัติการ}
\begin{enumerate}
    \item หากเลื่อนถาดอุปกรณ์ผ่าตัดให้ห่างจากหลอดรังสี UV-C เพิ่มขึ้นเป็น $3$ เท่าของระยะเดิม ความเข้มรังสี UV-C ที่ตกกระทบบนถาดจะลดลงเหลือกี่เท่า และต้องเปิดไฟนานขึ้นกี่เท่าจึงจะได้ปริมาณรังสีสะสม (UV Dose) เท่าเดิม
    \item เพราะเหตุใดรังสี UV-C จึงเหมาะสำหรับการฆ่าเชื้อเฉพาะบนพื้นผิวโต๊ะและอากาศ แต่ไม่สามารถใช้ฆ่าเชื้อในอาหารเลี้ยงเชื้อที่บรรจุอยู่ในขวดแก้วหนาหรือกล่องพลาสติกทึบได้
\end{enumerate}
